\section*{第 \ref{ch:g11:organic-skeletons} 章——有机分子：骨架与命名}

\begin{solution}{exo:g11:organic-skeletons:1}
两者都含六个碳\omterm{def:g7:atoms-and-molecules:atom}{原子}（端点与拐角），\omterm{def:g11:organic-skeletons:formulas}{分子式}均为 \ce{C6H14}：己烷与
2-甲基戊烷互为\omterm{def:g11:organic-skeletons:isomer}{同分异构体}。
\end{solution}

\begin{solution}{exo:g11:organic-skeletons:2}
丙烷；己烷；2-甲基丁烷。
\end{solution}

\begin{solution}{exo:g11:organic-skeletons:3}
\ce{CH3-CH2-CH3}; \chemfig{CH_3-CH(-[2]CH_3)-CH_3};
\chemfig{H-C(-[2]H)(-[6]H)-C(-[2]H)(-[6]H)-H}.
\end{solution}

\begin{solution}{exo:g11:organic-skeletons:4}
第一种为 3-甲基戊烷：含 6 个碳\omterm{def:g7:atoms-and-molecules:atom}{原子}；其三个 \ce{CH3} 端基
共含 9 个氢\omterm{def:g7:atoms-and-molecules:atom}{原子}，两个 \ce{CH2} 拐角含 4 个氢\omterm{def:g7:atoms-and-molecules:atom}{原子}，其
\ce{CH} 支化点含 1 个氢\omterm{def:g7:atoms-and-molecules:atom}{原子}：\ce{C6H14}。环己烷有 6 个拐角，每个都是
\ce{CH2}: \ce{C6H12}.
\end{solution}

\begin{solution}{exo:g11:organic-skeletons:5}
\ce{C8H18}。由 $2n + 2 = 20$ 得 $n = 9$：壬烷，\ce{C9H20}。
\end{solution}

\begin{solution}{exo:g11:organic-skeletons:6}
\setchemfig{atom sep=1.6em}
3-甲基己烷 \chemfig{-[:30]-[:-30]-[:30](-[:90])-[:-30]-[:30]-[:-30]}，
\ce{C7H16}；2,2-二甲基丁烷 \chemfig{-[:0](-[:90])(-[:-90])-[:0]-[:30]}，
\ce{C6H14}；3-乙基戊烷
\chemfig{-[:30]-[:-30](-[:-90]-[:-30])-[:30]-[:-30]}, \ce{C7H16}.
\end{solution}

\begin{solution}{exo:g11:organic-skeletons:7}
戊烷、2-甲基丁烷和 2,2-二甲基丙烷，
图示见本章图形。
\end{solution}

\begin{solution}{exo:g11:organic-skeletons:8}
“2-乙基丁烷”的最长链经过乙基，共五个碳\omterm{def:g7:atoms-and-molecules:atom}{原子}，
其正确名称为 3-甲基戊烷。“4-甲基戊烷”的编号从错误的一端开始，
正确应为 2-甲基戊烷。“1-甲基丁烷”实际上就是五个碳\omterm{def:g7:atoms-and-molecules:atom}{原子}的链，
即戊烷。
\end{solution}

\begin{solution}{exo:g11:organic-skeletons:9}
2,2-二甲基丙烷（\qty{9.5}{\celsius}）<2-甲基丁烷
（\qty{27.8}{\celsius}）<戊烷（\qty{36.1}{\celsius}）。同样的\omterm{def:g7:atoms-and-molecules:atom}{原子}
排列得越来越紧凑：\omterm{def:g7:atoms-and-molecules:molecule}{分子}间接触减少，
范德华相互作用减弱，沸点降低。
\end{solution}

\begin{solution}{exo:g11:organic-skeletons:10}
将链首尾相连成环，多出一个 \ce{C-C} 键，它占据了
两个 \ce{C-H} 键（每端各一个）的位置。\omterm{def:g11:organic-skeletons:formulas}{分子式}不同，
故不是\omterm{def:g11:organic-skeletons:isomer}{同分异构体}。环己烷不是\omterm{def:g11:organic-skeletons:alkane}{烷烃}（其骨架为环状，
\omterm{def:g11:organic-skeletons:formulas}{分子式}不符合 \ce{C_{n}H_{2n+2}}）。
\end{solution}

\begin{solution}{exo:g11:organic-skeletons:11}
丁烷与 2-甲基丙烷（\ce{C4H10}）。丙烷为 \ce{C3H8}，
环丁烷为 \ce{C4H8}。
\end{solution}

\begin{solution}{exo:g11:organic-skeletons:12}
\setchemfig{atom sep=1.5em}
\begin{center}
\small
\begin{tabular}{ccc}
\chemfig{-[:30]-[:-30]-[:30]-[:-30]-[:30]} &
\chemfig{-[:30](-[:90])-[:-30]-[:30]-[:-30]} &
\chemfig{-[:30]-[:-30](-[:-90])-[:30]-[:-30]} \\
己烷 & 2-甲基戊烷 & 3-甲基戊烷 \\[6pt]
\chemfig{-[:0](-[:90])(-[:-90])-[:0]-[:30]} &
\chemfig{-[:30](-[:90])-[:-30](-[:-90])-[:30]} & \\
2,2-二甲基丁烷 & 2,3-二甲基丁烷 &
\end{tabular}
\end{center}
\end{solution}

\begin{solution}{exo:g11:organic-skeletons:13}
两条长锯齿链可沿全长并排放置，接触的\omterm{def:g7:atoms-and-molecules:atom}{原子}很多；两个球只在
很小的面积上接触。范德华相互作用存在于相互接触的\omterm{def:g7:atoms-and-molecules:atom}{原子}之间：
直链\omterm{def:g7:atoms-and-molecules:molecule}{分子}之间的作用更大，
这些\omterm{def:g7:atoms-and-molecules:molecule}{分子}需要更高的温度才能分开。
\end{solution}

\begin{solution}{exo:g11:organic-skeletons:14}
最长链含六个碳\omterm{def:g7:atoms-and-molecules:atom}{原子}；从左端编号，取代基位于 2、4、4，
从右端编号为 3、3、5；首个不同的数字（2 与 3）决定编号，
故为 2,4,4-三甲基己烷。
\end{solution}

\begin{solution}{exo:g11:organic-skeletons:15}
\setchemfig{atom sep=1.6em}
\chemfig{-[:0](-[:90])(-[:-90])-[:-30]-[:30](-[:90])-[:-30]}：一条
含五个碳\omterm{def:g7:atoms-and-molecules:atom}{原子}的链，碳 2 上连两个甲基，碳 4 上连一个甲基，
共八个碳\omterm{def:g7:atoms-and-molecules:atom}{原子}；$2 \times 8 + 2 = 18$，氢\omterm{def:g7:atoms-and-molecules:atom}{原子}数为 18，
\ce{C8H18}，为辛烷的\omterm{def:g11:organic-skeletons:formulas}{分子式}，二者互为\omterm{def:g11:organic-skeletons:isomer}{同分异构体}。
\omterm{def:g11:organic-skeletons:formulas}{结构简式}：\chemfig{CH_3-C(-[2]CH_3)(-[6]CH_3)-CH_2-CH(-[2]CH_3)-CH_3}。
\end{solution}

\begin{solution}{pb:g11:organic-skeletons:1}
\setchemfig{atom sep=1.5em}
\textbf{1.} \ce{C4H10}：当 $n = 4$ 时，$2n + 2 = 10$。

\textbf{2.} \chemfig{H-C(-[2]H)(-[6]H)-C(-[2]H)(-[6]H)-C(-[2]H)(-[6]H)-C(-[2]H)(-[6]H)-H}.

\textbf{3.} \ce{CH3-CH2-CH2-CH3}; \chemfig{-[:30]-[:-30]-[:30]}.

\textbf{4.} \chemfig{H-C(-[2]H)(-[6]H)-C(-[6]H)(-[2,1.5]C(-[1]H)(-[3]H)(-[2]H))-C(-[2]H)(-[6]H)-H};
\chemfig{CH_3-CH(-[2]CH_3)-CH_3}; \chemfig{-[:30](-[:90])-[:-30]}。它与前者
\omterm{def:g7:atoms-and-molecules:atom}{原子}相同，都是四个碳、十个氢，
但成键次序不同。

\textbf{5.} 是：\omterm{def:g11:organic-skeletons:isomer}{构造异构体}（碳骨架不同）。

\textbf{6.} 三者的沸点均低于
\qty{20}{\celsius}。

\textbf{7.} 在 \qty{-5}{\celsius} 时，低于其沸点，丁烷
在常压下仍为液体，几乎不逸出气体。

\textbf{8.} 丙烷（\qty{-42}{\celsius}）和 2-甲基丙烷
（\qty{-11.7}{\celsius}）的沸点均低于冬季气温，它们
在寒冷天气中仍能提供大量气体。

\textbf{9.} 2-甲基丙烷带支链，结构更紧凑，\omterm{def:g11:polarity-and-cohesion:van-der-waals}{范德瓦
耳斯相互作用}较弱。

\textbf{10.} 含三个碳\omterm{def:g7:atoms-and-molecules:atom}{原子}时，唯一可能的骨架为
链 \ce{C-C-C}：任何碳\omterm{def:g7:atoms-and-molecules:atom}{原子}移到别处，得到的仍是同一条链。

\textbf{11.} 2-甲基丁烷。

\textbf{12.} 主链应从靠近取代基的一端编号，
故取代基的位次为 2，而非 3。

\textbf{13.} 2、3 和 5。

\textbf{14.} 庚烷，\chemfig{-[:30]-[:-30]-[:30]-[:-30]-[:30]-[:-30]}。

\textbf{15.} 2-甲基己烷 \chemfig{-[:30](-[:90])-[:-30]-[:30]-[:-30]-[:30]}
和 3-甲基己烷 \chemfig{-[:30]-[:-30](-[:-90])-[:30]-[:-30]-[:30]}。
“4-甲基己烷”是从错误一端编号的 3-甲基己烷。

\textbf{16.} 2,2-二甲基戊烷、2,3-二甲基戊烷、
2,4-二甲基戊烷和 3,3-二甲基戊烷。

\textbf{17.} 是，3-乙基戊烷。若乙基连在碳 2 上，
则会通过它形成六碳主链，该\omterm{def:g7:atoms-and-molecules:molecule}{分子}实为 3-甲基己烷。

\textbf{18.} 2,2,3-三甲基丁烷，
\chemfig{-[:0](-[:90])(-[:-90])-[:0](-[:90])-[:0]}。

\textbf{19.} $1 + 2 + 4 + 1 + 1 = 9$。

\textbf{20.} 庚烷共有 9 种\omterm{def:g11:organic-skeletons:isomer}{构造异构体}（包括其本身）。
\end{solution}
